\chapter{Documentação, reprodutibilidade e uso de IA}

\section{Organização do projeto}

A versão atual do repositório contém código-fonte, testes, exemplos, documentação, declaração de IA, registro de contribuições e diagramas. Cada ER possui um arquivo \texttt{.jff} para JFLAP e um arquivo \texttt{.md} com Mermaid. Os seis arquivos \texttt{.jff} foram carregados e simulados pelo próprio JFLAP 7.1 (\texttt{tests/validate\_jflap\_runtime.py}) em 40.044 casos, sem divergências.

\section{Reprodutibilidade}

A execução dos testes é realizada com:

\begin{lstlisting}[language=bash]
pip install -r requirements.txt
pytest -q
python tests/validate_all.py
python tests/validate_afne_structure.py
\end{lstlisting}

As classes utilizadas nos diagramas JFLAP são explicitadas, por exemplo, como \texttt{[A-Za-z]}, \texttt{[0-9]}, \texttt{[1-9]}, \texttt{[0-5]} e \texttt{[0-3]}, enquanto o simulador Python expande os conjuntos para seus símbolos constituintes.

\section{Uso de Inteligência Artificial}

Ferramentas de IA foram utilizadas como apoio na organização e revisão da especificação, análise das ERs e AFN$\varepsilon$, implementação e revisão do simulador e do lexer, criação dos testes, conferência dos diagramas e revisão documental. As ferramentas utilizadas foram o Claude (Anthropic) e o ChatGPT (OpenAI). O projeto mantém uma declaração explícita em \texttt{DECLARACAO\_IA.md}. A responsabilidade pela compreensão, revisão e alteração do material permanece com a equipe.

\section{Registro de contribuições}

As contribuições de cada integrante estão registradas em \texttt{CONTRIBUICOES.md}:

\begin{itemize}
\item \textbf{Igor Cecim Vilhena}: definição e formalização das Expressões Regulares (alfabetos, linguagens, notação formal e sintaxe Python), redação do relatório técnico e elaboração dos slides da apresentação.
\item \textbf{Andrey Lourival Andrade Garcia}: especificação da MiniLang, construção dos AFN$\varepsilon$ de Thompson e do simulador, implementação do MiniLexer, testes automatizados e validação AFN$\varepsilon$ $\times$ regex, diagramas JFLAP e Mermaid, exemplos, README, fichas explicativas das ERs e roteiro de comandos da apresentação.
\end{itemize}
